\section{Evolution of the Project Scope}
\label{sec:evolution}

The scope of this project has evolved significantly over its lifecycle, moving from broad conceptual aspirations to a highly focused, empirical engineering evaluation. This evolution was guided by continuous literature review, feasibility assessments, and early prototype experiments.

The initial conception of the project (as outlined in Proposal 1) began with a broad motivation to tackle the blockchain scalability trilemma. The early goals focused on conceptualizing a comprehensive ZK-rollup solution that could handle massive throughput while maintaining robust security properties. However, during the initial literature review phase, it became evident that simply building another general-purpose rollup would not contribute uniquely to the academic field.

Subsequently, as detailed in Proposal 2, the topic was narrowed significantly. The project pivoted away from building a production-grade system and instead adopted the objective of creating a \textit{modular benchmarking framework}. The focus shifted to empirically evaluating specific architectural trade-offs within ZK-rollups, specifically observing the interplay between sequencer scheduling, batch sizes, and data availability. The refined problem statement formalized the necessity of a controlled experimental environment capable of isolating performance bottlenecks under synthetic workloads.

As development progressed, the implementation phase (documented in the Progress Report) brought further practical refinements. The engineering team successfully implemented the core modular components—a sequencer, an executor, a submitter, and smart contracts—using Rust and Solidity. A critical realization during this phase was the immense computational overhead required for real cryptographic proof generation. To maintain testing velocity and isolate queueing dynamics from hardware limitations, the team introduced a dual-prover architecture. This architecture utilized a mock prover with parameterized delays for broad scheduling experiments, while restricting the real RISC0 zkVM prover to focused, cycle-analysis experiments.

This final report reflects the culmination of this evolution. It presents an accurate, empirical evaluation of batching logic, sequencing heuristics, proving overhead, and data availability strategies. It strictly separates the implemented research prototype's validated findings from the theoretical performance ceilings of production zkEVMs. By embracing a focused, transfer-centric execution engine, the framework successfully isolates and quantifies the critical throughput--latency--cost trade-offs that govern Layer-2 scalability.